%%%%%%%%%%%%%%%%%%%%%%%%%%%%%%%%%%%%%%%%%%%%%%%%%%%%%%%%%%%
\chapter*{Apport du projet}
\addcontentsline{toc}{chapter}{Apport du projet}
\label{ch:apport}
%%%%%%%%%%%%%%%%%%%%%%%%%%%%%%%%%%%%%%%%%%%%%%%%%%%%%%%%%%%

Ce Projet de Fin d'Études a été un catalyseur de transformation personnelle et professionnelle, nous confrontant à la dualité entre l'idéal théorique et le pragmatisme technique du terrain.

\paragraph{Une démarche de résolution de problèmes :}
Au-delà de la technique, ce travail nous a enseigné qu'innover ne se résume pas à empiler des outils, mais à structurer une démarche scientifique. Chaque étape, de l'extraction des besoins à l'inférence des modèles d'IA, a agi comme une leçon de résilience. Nous avons appris à valider nos hypothèses par l'expérimentation constante, transformant chaque obstacle en une opportunité d'affiner notre compréhension du système.

\paragraph{Polyvalence technologique et méthodologique :}
L'acquisition de compétences a été multidimensionnelle :
\begin{itemize}[noitemsep, topsep=1pt]
    \item \textbf{Agilité (Scrum) :} L'adoption du cadre Scrum nous a appris l'importance de la livraison incrémentale et de la flexibilité face aux imprévus.
    \item \textbf{Rigueur rédactionnelle (\LaTeX) :} L'utilisation de \LaTeX a imposé une discipline de fer dans la structuration et la présentation de la documentation technique.
    \item \textbf{Architecture Microservices :} Naviguer entre Java, Python et JavaScript nous a permis de maîtriser la communication au sein d'un écosystème hétérogène.
\end{itemize}

\paragraph{Défis rencontrés et résilience :}
Le parcours a été jalonné d'embûches. Les délais stricts et les limitations de puissance de calcul ont parfois freiné nos ambitions. L'intégration de modèles \textbf{Deep Learning} gourmands sur un environnement standard a nécessité une optimisation constante pour garantir la viabilité du prototype sans sacrifier ses performances.

\paragraph{Analyse critique et perspectives :}
Avec le recul, une planification plus granulaire aurait permis de mieux anticiper les risques de latence. Pour une itération future, nous privilégierions une infrastructure Cloud optimisée GPU pour accélérer les traitements de vision. Une documentation interne plus exhaustive dès le départ aurait également facilité la maintenance.

En conclusion, ce projet a été le pont entre notre parcours académique et nos ambitions professionnelles. Il nous a doté d'une vision systémique essentielle pour notre future carrière d'ingénieur.
